\documentclass[11pt,a4paper]{article}
\usepackage[utf8]{inputenc}
\usepackage[T1]{fontenc}
\usepackage{amsmath,amssymb}
\usepackage[margin=2.5cm]{geometry}
\usepackage{xcolor}

\setlength{\parindent}{0pt}
\setlength{\parskip}{6pt}

\begin{document}

\begin{center}
{\large Technical Neural Networks Assignment Sheet 2b}\\[4pt]
{\bfseries Names: Mohamadhosein Sabaghi (50547423)}
\end{center}

{\small (Programming Assignment PA-A is solved separately.)}

\section*{Assignment 7}

The $\delta$-rule as given in the lecture:
\[
  {}^p\!\Delta w_{gh} = \eta \cdot {}^p\delta_h \cdot {}^p\text{out}_g
\]

\begin{itemize}
  \item ${}^p\Delta w_{gh}$: the change applied to the weight from neuron $g$ to neuron $h$, for pattern $p$.
  \item $\eta$: the learning rate, controls the step size.
  \item ${}^p\delta_h$: the error signal of neuron $h$, defined as $\delta_h = -\partial E/\partial \mathrm{net}_h$. For an output
  neuron $m$: $\delta_m=(\hat y_m-y_m)\,f'(\mathrm{net}_m)$. For a hidden neuron $h$:
  $\delta_h = f'(\mathrm{net}_h)\sum_k \delta_k w_{hk}$, i.e.\ the $\delta$'s of the neurons it feeds into, propagated
  backward through the weights.
  \item ${}^p\text{out}_g$: the output (activation) of the pre-synaptic neuron $g$, i.e.\ the value that travels along
  this connection.
\end{itemize}

\section*{Assignment 8}

Weights are counted including the BIAS, so between a layer of $N$ inputs and a layer of $H$ neurons there are
$(N+1)\cdot H$ weights.

Two hidden layers $N$-$H_1$-$H_2$-$M$:
\[
  W_2 = (N+1)H_1 + (H_1+1)H_2 + (H_2+1)M
\]

One hidden layer $N$-$H$-$M$:
\[
  W_1 = (N+1)H + (H+1)M = H(N+1+M) + M
\]

Setting $W_1=W_2$ and cancelling $M$ on both sides:
\[
  H(N+1+M) = (N+1)H_1 + H_1H_2 + H_2 + H_2M = (N+1)H_1 + H_2(H_1+1+M)
\]
\[
  \boxed{H = \dfrac{(N+1)H_1 + H_2(H_1+1+M)}{N+1+M}}
\]

For $N$-$H_1$-$H_2$-$M = 5$-$21$-$30$-$4$:
\[
  H = \frac{6\cdot21 + 30\cdot(21+1+4)}{5+1+4} = \frac{126+780}{10} = \frac{906}{10} = 90.6
\]

$H$ must be an integer, so we round to $H=91$. Check: original network has
$W_2 = 6\cdot21+22\cdot30+31\cdot4 = 126+660+124 = 910$ weights; the replacement $5$-$91$-$4$ network has
$W_1 = 6\cdot91+92\cdot4 = 546+368 = 914$ weights -- very close to the original 910.

\section*{Assignment 9}

Assume a simple Perceptron (inputs $x_1,x_2$, BIAS, weights $w_1,w_2,w_b$, step function) could realize XOR:
\[
  y = 1 \text{ if } \mathrm{net}=w_1x_1+w_2x_2+w_b>0, \qquad y=0 \text{ otherwise.}
\]

XOR requires $y(0,0)=0,\ y(0,1)=1,\ y(1,0)=1,\ y(1,1)=0$, which gives four inequalities:
\[
  (0,0):\ w_b \le 0 \quad (\mathrm{I}) \qquad\qquad
  (0,1):\ w_2+w_b > 0 \quad (\mathrm{II})
\]
\[
  (1,0):\ w_1+w_b > 0 \quad (\mathrm{III}) \qquad\qquad
  (1,1):\ w_1+w_2+w_b \le 0 \quad (\mathrm{IV})
\]

Adding (II) and (III):
\[
  w_1+w_2+2w_b > 0 \;\;\Longrightarrow\;\; w_1+w_2 > -2w_b \quad (\mathrm{V})
\]
From (IV): $w_1+w_2 \le -w_b$. Combined with (V):
\[
  -2w_b < w_1+w_2 \le -w_b \;\;\Longrightarrow\;\; -2w_b < -w_b \;\;\Longrightarrow\;\; w_b > 0
\]

This contradicts (I), which requires $w_b\le 0$. No weights $w_1,w_2,w_b$ can satisfy all four inequalities at once,
so a simple Perceptron without a hidden layer cannot implement XOR. $\blacksquare$

\section*{Assignment 10b}

\[
  {}^pE^*(w_{hm}) = \frac12\sum_{m=0}^{M}\left({}^p\hat y_m - {}^py_m\right)^6
\]

As in the lecture, we use gradient descent, $\Delta w_{hm}=-\eta\,\partial E^*/\partial w_{hm}$, and the chain rule
through $\mathrm{net}_m$:
\[
  \frac{\partial E^*}{\partial w_{hm}} = \frac{\partial E^*}{\partial y_m}\cdot\frac{\partial y_m}{\partial \mathrm{net}_m}\cdot\frac{\partial \mathrm{net}_m}{\partial w_{hm}}
\]

Each factor:
\[
  \frac{\partial E^*}{\partial y_m} = \frac12\cdot 6(\hat y_m-y_m)^5\cdot(-1) = -3(\hat y_m-y_m)^5,
  \qquad
  \frac{\partial y_m}{\partial \mathrm{net}_m} = f'(\mathrm{net}_m),
  \qquad
  \frac{\partial \mathrm{net}_m}{\partial w_{hm}} = \mathrm{out}_h
\]

So:
\[
  \frac{\partial E^*}{\partial w_{hm}} = -3(\hat y_m-y_m)^5\,f'(\mathrm{net}_m)\,\mathrm{out}_h
\]

As in the lecture, define $\delta_m^* := -\partial E^*/\partial \mathrm{net}_m$, so that the update keeps the usual
$\Delta w_{hm}=\eta\,\delta_m^*\,\mathrm{out}_h$ form:
\[
  \boxed{{}^p\delta_m^* = 3\left({}^p\hat y_m - {}^py_m\right)^5 f'({}^p\mathrm{net}_m)}
\]
\[
  {}^p\Delta w_{hm} = \eta\,{}^p\delta_m^*\,{}^p\mathrm{out}_h = 3\eta\left({}^p\hat y_m-{}^py_m\right)^5 f'({}^p\mathrm{net}_m)\,{}^p\mathrm{out}_h
\]

\section*{Assignment 11b}

\textbf{AI use:} I used Claude (Anthropic) to work through this part, as the sheet requires. I gave it the
assignment text directly and asked for the derivative calculations for $\tanh$ and the logistic function, shown
step by step. The derivation worked on the first try and I checked each step by hand, so nothing was discarded.

\textbf{Logistic function.} $f(z)=(1+e^{-z})^{-1}$. By the chain rule:
\[
  f'(z) = -(1+e^{-z})^{-2}\cdot(-e^{-z}) = \frac{e^{-z}}{(1+e^{-z})^2}
\]
Since $1-f(z) = 1-\dfrac{1}{1+e^{-z}} = \dfrac{e^{-z}}{1+e^{-z}}$, we get
\[
  f(z)\,(1-f(z)) = \frac{1}{1+e^{-z}}\cdot\frac{e^{-z}}{1+e^{-z}} = \frac{e^{-z}}{(1+e^{-z})^2} = f'(z)
\]
\[
  \boxed{f'(z) = f(z)\,(1-f(z))}
\]

\textbf{tanh.} With $u=e^z-e^{-z}$, $v=e^z+e^{-z}$, so $\tanh(z)=u/v$, $u'=v$, $v'=u$. By the quotient rule:
\[
  \tanh'(z) = \frac{u'v-uv'}{v^2} = \frac{v^2-u^2}{v^2} = 1-\left(\frac{u}{v}\right)^2
\]
\[
  \boxed{\tanh'(z) = 1-\tanh^2(z)}
\]

\section*{Assignment 12b}

Claim: $\tanh(z) = 2f_{\log}(2z)-1$, where $f_{\log}(z)=\dfrac{1}{1+e^{-z}}$.

\[
  2f_{\log}(2z)-1 = \frac{2}{1+e^{-2z}}-1 = \frac{2-(1+e^{-2z})}{1+e^{-2z}} = \frac{1-e^{-2z}}{1+e^{-2z}}
\]

Multiplying numerator and denominator by $e^{z}$:
\[
  \frac{1-e^{-2z}}{1+e^{-2z}} = \frac{e^{z}-e^{-z}}{e^{z}+e^{-z}} = \tanh(z)
\]

So $\tanh(z) = 2f_{\log}(2z)-1$: $\tanh$ is exactly the logistic function, stretched horizontally by a factor $2$
and then shifted/rescaled from the range $(0,1)$ to $(-1,1)$. $\blacksquare$

\end{document}
